\subsection{Nonunit tensor powers and the actual level normalization}

\begin{theorem}[DP50: the exact effect of a $p$-divisible tensor exponent]\label{DP50}
Use the arbitrary-unit singleton/complement display and its actual primitive Poincar\'e kernel from DP46. Assume $p$ odd, $h\ge3$, $q=p^h$, $d=h-1$, and $R=W(k)$ with perfect $k\supset\F_q$. Put
\[
 s=\frac1{p-1},\qquad \rho=\frac1{q-1},\qquad
 r_i=s-p^i\rho\quad(1\le i\le d).
\]
Let $u\in R^\times$, let $\mathbf v\in R^d$ be primitive, and set
\[
 j=\min\{i:v_i\in R^\times\},\qquad
 k_*=h-j,\qquad S_{k_*}=\frac{p^{k_*}-1}{p-1}.
\]
For $\kappa=p^a\lambda$ with $a\ge1$ and $\lambda\in R^\times$, write
\[
 K_\kappa(t)=\exp\bigl(\kappa\beta(ut,\mathbf vt)\bigr)
       =\sum_{n\ge0}A_nt^n.
\]
Then, for every rational $w\ge0$,
\begin{equation}\label{eq:nonunitradius}
 \nu_w(K_\kappa)=
 \begin{cases}-\infty,&w<r_j,\\0,&w\ge r_j.\end{cases}
\end{equation}
More precisely, for every $m\ge0$,
\begin{equation}\label{eq:nonunitsequence}
 N_m=p^{k_*}q^m+1,\qquad
 \vp(A_{N_m})=a-S_{k_*}-r_jp^{k_*}(q^m-1).
\end{equation}
Consequently
\begin{equation}\label{eq:nonunitslope}
 \liminf_{n\to\infty}\frac{\vp(n!A_n)}n
       =s-r_j=\frac{p^j}{q-1}.
\end{equation}
These statements, including the exact valuations in
\eqref{eq:nonunitsequence}, hold in every normalized invariant integral frame of the same primitive biextension.

For every integral formal function $U$, put
$V(t)=U(E_G(ut),E_H(\mathbf vt))$ and
$V(t)K_\kappa(t)=\sum_n C_nt^n$. Then
\begin{equation}\label{eq:nonunitprefactor}
 \vp(n!C_n)\ge\frac{p^j}{q-1}n-\frac{s_p(n)}{p-1}.
\end{equation}
This is a sharp slope as a bound valid for all integral $U$, because $U=1$ attains \eqref{eq:nonunitslope}. Equality of asymptotic slopes is not asserted for each individual prefactor.
\end{theorem}

\begin{proof}
\emph{The boundary contraction.}
Set
\[
 \Delta=r_j-\rho>0,\qquad
 \omega_0=r_j,\quad\omega_i=r_i\ (1\le i\le d).
\]
The smallest possible numerator of
\[
 \Delta=\frac{1+p+\cdots+p^{h-1}-p^j-1}{q-1}
\]
occurs at $j=d$ and is $p+\cdots+p^{d-1}>0$; this uses $h\ge3$.
Choose $e$ clearing the weight denominators, adjoin $\pi$ with $\pi^e=p$, and put
$\alpha_i=\pi^{e\omega_i}$ and $\gamma=\alpha_0=\alpha_j$.
The contracted scalar logarithm reduces to the identity, since its nonlinear coefficient at degree $q^n$ has valuation
$(q^n-1)r_j-n>0$.

For the dual, every selected first arrow has weight defect one. Its ordinary arrows have degree $p$ and its exceptional arrow degree $p^2$.
A path through $b\ge2$ selected arrows has weight increment at least
$1+p+\cdots+p^{b-1}>b$; its denominator $p^b$ therefore leaves positive valuation.
DP40's exact path logarithm, including arbitrary unit coefficients, gives
\begin{equation}\label{eq:nonunitreducedlog}
 \overline{\widetilde L}_G=S,\qquad
 \overline{\widetilde L}_H=\id+M_\eta,
\end{equation}
where
\[
 (M_\eta)_1=\eta_1Y_d^{p^2},\qquad
 (M_\eta)_i=\eta_iY_{i-1}^p\ (2\le i\le d),
 \qquad\eta_i\in k^\times.
\]
The scaled logarithms and their formal inverses are integral with identity linear part. Since $M_\eta$ is additive over the residue field,
\[
 \overline{\widetilde E}_H
       =\sum_{b\ge0}(-1)^bM_\eta^{\circ b}
\]
in the formal degree topology. Its coefficients need not be Frobenius-fixed.

The contracted input to dual coordinate $i$ is $v_i\gamma/\alpha_i$.
For $i>j$ its valuation is positive because $r_i<r_j$.
For $i<j$, use $v_i\in pR$ and
$0<r_i-r_j=(p^j-p^i)\rho<1$ to reach the same conclusion.
Exactly the $j$th unit input survives reduction.
The unique path from $j$ to $d$, followed by any number of complete cycles, therefore gives
\begin{equation}\label{eq:nonunitinversecycle}
 \overline{\widetilde E}_{H,d}(z(T))
       =\sum_{m\ge0}\xi_mT^{p^{k_*-1}q^m},
       \qquad\xi_m\in k^\times.
\end{equation}
The exponents are distinct and there is one path for each, so none of these coefficients cancels. The scalar inverse reduces to $\bar uT$.

\emph{The entire primitive logarithm at this boundary.}
DP46's variable-ghost formula is
\begin{align*}
 b(S,\mathbf Y)
 ={}&-\sum_{i=1}^d\sum_{\ell,n\ge0}
 \epsilon_{i,\ell,n}
 \frac{S^{p^{n+\ell h+i}}Y_i^{p^n}}{p^{n+\ell+1}}\\
 &+\left(\sum_{b\ge0}U_b\frac{S^{q^b}}{p^b}\right)
   \left(\sum_{i=1}^d\sum_{\ell\ge0}
       V_{i,\ell}\frac{Y_i^{p^{\ell h+h-i}}}{p^{d\ell+h-i}}\right),
\end{align*}
where the displayed coefficients are units and $U_0=1$.
Its $SY_d^p$ coefficient is $\upsilon_d^{-1}/p$, where
$\upsilon_d=w_0(u_d)\in R^\times$ is the zeroth ghost of the last display arrow.
Put
\[
 \tau=r_j+pr_d-1=s-1+\Delta,\qquad
 \mu=\alpha_0\alpha_d^p/p,
 \quad\vp(\mu)=\tau.
\]
We verify that every other displayed term has strictly larger contracted valuation, even though $\tau$ may be negative.

For the first family set $A_i^*=p^ir_j+r_i=s+p^i\Delta>s$.
Its baseline excess is
\[
 A_i^*-1-\tau=(p^i-1)\Delta>0.
\]
The additional excess at general $n,\ell$ is
\[
 (p^n-1)A_i^*-n+p^{n+i}r_j(q^\ell-1)-\ell.
\]
Both parts are nonnegative, strictly positive for a positive corresponding index: use $(p^n-1)s\ge n$, $A_i^*>s$, and $r_j(q-1)>1$.
For the product family put $e_i=h-i$. The baseline excess is
\[
 r_j+p^{e_i}r_i-e_i-\tau
       =(p^{e_i}-p)s-e_i+1\ge0,
\]
with equality exactly at $e_i=1$.
The extra excess at positive line index $b$ or dual cycle index $\ell$ is
\[
 r_j(q^b-1)-b+p^{e_i}r_i(q^\ell-1)-d\ell>0.
\]
For the last assertion use
\[
 p^{e_i}r_i(q-1)\ge\frac{q-p}{p-1}
       =p+p^2+\cdots+p^{h-1}>d.
\]
Thus the unique minimum is the $+\upsilon_d^{-1}SY_d^p/p$ term:
\begin{equation}\label{eq:nonunitloginitial}
 \overline{\mu^{-1}b(\alpha_0S,\alpha_1Y_1,\ldots,\alpha_dY_d)}
       =\bar\upsilon_d^{-1}SY_d^p.
\end{equation}
This checks the full ghost series; it does not extrapolate the narrower interval from the unit-exponent proof.

\emph{An infinite initial form survives the exponential.}
Compose \eqref{eq:nonunitloginitial} with the integral inverse maps.
By \eqref{eq:nonunitinversecycle},
\begin{equation}\label{eq:nonunitinfiniteinitial}
 \overline{\mu^{-1}\beta(u\gamma T,\mathbf v\gamma T)}
 =\bar\upsilon_d^{-1}\bar uT
   \left(\sum_{m\ge0}\xi_mT^{p^{k_*-1}q^m}\right)^p
 =\sum_{m\ge0}\chi_mT^{N_m},
 \qquad\chi_m\in k^\times.
\end{equation}
The restricted logarithm consequently has infimum $\tau$, attained at infinitely many coefficients after scaling.
For $z(T)=\kappa\beta(u\gamma T,\mathbf v\gamma T)$ its infimum is
\[
 L=a+\tau=a+s-1+\Delta>s.
\]
For every exponential order $b\ge2$,
\[
 \nu(z^b/b!)-L
       \ge(b-1)L-\vp(b!)
       \ge(b-1)(L-s)>0.
\]
It follows that $\exp z$ and its inverse are integral and that the entire initial form of $\exp z-1$, divided by $\kappa\mu$, is the same nonzero series as in \eqref{eq:nonunitinfiniteinitial}. Therefore
\[
 \vp(A_{N_m})+r_jN_m=a+\tau
       \qquad(m\ge0).
\]
The identity $\tau-r_j(p^{k_*}+1)=-S_{k_*}$ proves
\eqref{eq:nonunitsequence}. Its right side is integral, since
$r_j(q-1)=1+p+\cdots+p^{h-1}-p^j$ is an integer.

At weight $r_j$ the kernel has infimum zero, and the same holds at every larger weight. At a smaller weight $w$, its displayed coefficients satisfy
\[
 \vp(A_{N_m})+wN_m=L-(r_j-w)N_m\longrightarrow-\infty.
\]
This proves \eqref{eq:nonunitradius} and the coefficient slope $-r_j$ directly. Legendre's formula gives \eqref{eq:nonunitslope}; no radius assertion has been inferred from the logarithm alone.

\emph{Frames and prefactors.}
Let $J$ compare two normalized invariant integral frames.
Its axes are one. Every positive supported monomial has positive $S$ degree and, under the original weights $(\rho,r_1,\ldots,r_d)$, has weight in $s\Z_{>0}$ by the torus quantization of DP46.
Increasing the scalar weight by $\Delta$ gives
\[
 \nu(J-1)\ge s+\Delta>s,\qquad
 \nu(\log J)\ge s+\Delta.
\]
The logarithm estimate follows termwise from $\vp(b)\le(b-1)s$.
The invariant splitting changes by $\pm\log J$, with no invariant bilinear remainder. Its contribution after multiplication by $\kappa$ has valuation at least
$a+s+\Delta=L+1$, and cannot change the infinite initial form or its coefficient valuations.

Finally $V(\gamma T)$ is integral for every integral $U$, by the integral inverse maps and inputs already checked. Hence $V K_\kappa$ is integral at weight $r_j$ and
$\vp(C_n)\ge-r_jn$. This proves \eqref{eq:nonunitprefactor}.
The choice $U=1$ proves universal sharpness. An arbitrary prefactor is not multiplied by $\kappa$ as a frame logarithm is, so it can enter at the first nonconstant boundary weight. Equality for each such prefactor does not follow from this argument.
\end{proof}

Thus one factor of $p$ changes the pure-kernel factorial slope from
$1/((p-1)(p^{h-j}+1))$ to $p^j/(p^h-1)$; further factors only change the constant valuation in \eqref{eq:nonunitsequence}.
For the actual height-three block at $p=11$ of DP49, these values are
\[
\begin{array}{c|c|c}
\text{dual residue class}&p\nmid\kappa&p\mid\kappa\\\hline
v_1\in R^\times&1/1220&11/1330\\
v_1\in11R,\ v_2\in R^\times&1/120&121/1330
\end{array}
\]
The last column is exact for the pure kernel and the sharp universal lower-bound slope with integral prefactors. These are genuine local improvements, not the result of adding $\vp(\kappa)$ to an old asymptotic slope.

\begin{corollary}[DP50: the recovered level denominator outweighs this improvement]
The recovered definition \path{thExpr} in
\path{Formalism/LocalToGlobal/ProjectiveLattices/IsThRule.lean}, lines 66--70, uses tensor exponent $M^2$ at lines 68 and 70 and the factor $M^{-P.\mathrm{totalDegree}}$ at line 69.
For a homogeneous degree-$n$ ordinary-path specialization, let $b=\vp(M)\ge1$.
Then the normalized pure kernel $K_{M^2}(t/M)$ has exact factorial slope
\[
 \liminf_n\frac{\vp(n!A_n/M^n)}n
       =\frac{p^j}{q-1}-b<0.
\]
Relative to a unit-level pure kernel, the change is
\[
 \frac{p^j}{q-1}-b-\frac{s}{p^{h-j}+1}
       <s-b\le-\frac b2.
\]
Thus forcing $p\mid M$ cannot improve this universal ordinary-path slope after the actual normalization.
\end{corollary}

\begin{proof}
The tensor exponent $M^2$ has valuation $2b\ge1$, so DP50 applies.
Every coefficient of $K_{M^2}(t/M)$ is $A_n/M^n$, whose valuation is decreased by $bn$.
This proves the exact slope.
Since $p^j/(q-1)<s\le1/2$ and $b\ge1$, subtracting the positive unit-level slope gives the displayed comparison.
For an arbitrary integral prefactor the corresponding sharp universal lower bound is
\[
 \vp(n!C_n/M^n)
       \ge\left(\frac{p^j}{q-1}-b\right)n
           -\frac{s_p(n)}{p-1}.
\]
Its universal sharpness again uses the pure kernel; no equality for every individual section has been asserted.
\end{proof}

This corollary concerns the actual displayed tensor/degree normalization and a homogeneous ordinary-path specialization. It allows additional divisibility from a particular algebraic section, a special selected harmonic value, or a compensating factor elsewhere in the construction; none is supplied merely by $p\mid M$.
For a fixed nonzero integer $\kappa$, its exceptional odd primes also satisfy the elementary bound
\[
 0\le\sum_{\substack{p\mid\kappa\\p\ \mathrm{odd}}}\log p
 \left(\frac{p^{j_p}}{p^{h_p}-1}
       -\frac1{(p-1)(p^{h_p-j_p}+1)}\right)
 \le\sum_{\substack{p\mid\kappa\\p\ \mathrm{odd}}}\frac{\log p}{p-1}
 \le\frac12\log|\kappa|.
\]
The difference in parentheses is positive: its numerator after a common positive denominator is
$(p-2)p^{h_p}+(p-1)p^{j_p}+1$.
At derivative order $n$, these finitely many primes contribute only $O_\kappa(n)$ to a prime-sum comparison, without changing an $n\log X$ coefficient as the cutoff $X$ grows. This does not dismiss a finite improvement in a finite inequality, or supply a cost theorem when $\kappa$ varies with the construction.

A rational quadratic gauge remains a separate input. The new radius $r_j$ can be below $s/2$, so an arbitrary integral quadratic coefficient is not automatically harmless as it was in the unit-exponent situation of DP48. If the quadratic is $\kappa Q$ with integral $Q$, its extra factor of $p$ makes its exponential and inverse integral at every nonnegative weight. No selected harmonic determinant or shorter global chain is identified by the ordinary-path theorem alone.

